\bookchapter{Заключение: Сохраняйте любопытство, сохраняйте ясность}

Поздравляем, что дочитали до этого места.

Эта книга прошла путь от «железа» телефона до программной экосистемы, от повседневных настроек до режима разработчика, от защиты данных до цифрового наследия. Если она помогла вам решить хотя бы одну практическую задачу — пусть даже научиться выключать лишнее разрешение или завести привычку делать резервные копии — значит, эти строки написаны не зря.

Но важнее любых конкретных приёмов — две вещи:

\begin{itemize}
  \item \textbf{Сохраняйте любопытство}: телефон меняется каждый год, меняется ИИ, меняются складные экраны, меняется спутниковая связь. Операции, которые вы выучите сегодня, завтра могут устареть, но любопытство «хочу знать почему» не устареет никогда. Лишний вопрос «зачем нужно это разрешение», «где хранятся эти данные», «не раскроет ли эта функция конфиденциальность» полезнее, чем запоминание сотни операций.
  \item \textbf{Сохраняйте ясность}: телефон — это инструмент, а не хозяин. Он спроектирован так, чтобы привлекать ваше внимание, занимать ваше время и направлять ваше потребление. Использовать его — ваш выбор; быть использованным им — ваша потеря контроля. Периодически смотрите на экранное время, периодически задавайте себе вопрос: «я пользуюсь телефоном для дела или телефон «стачивает» меня?» — это не позёрство, это инстинкт выживания в цифровую эпоху.
\end{itemize}

Технологии меняются, но человеческие потребности неизменны — нам нужна безопасность, нужны связи, нужно уважение. Телефон — лишь носитель. Настоящий «курс по смартфону» не в книге, а в моменте, когда вы каждый день разблокируете экран.

\begin{tcbpractice}{Последнее упражнение}
Закройте эту книгу, возьмите телефон и сделайте одну вещь:\textbf{за 30 секунд подумайте, какие три вещи в вашем телефоне «по-настоящему было бы жалко потерять»?}Затем убедитесь, что у них есть резервная копия. И книга свою задачу выполнила.

\end{tcbpractice}
